
Historická šifrácia je fascinujúca oblasť skúmania nielen preto, že odhaľuje metódy utajovania informácií v minulosti, ale aj preto, že poskytuje zaujímavý úvod do politiky, vojny a spoločenských pomerov jej podaktorých častí. Medzi pokročilé a historicky významné šifrovacie systémy patrili takzvané nomenklátory. Tieto systémy párujúce rôzne techniky zašifrovávania jednotlivých písmen ako sú substitúcie, n-gramy, klamáče (písmená alebo iné symboly, ktoré sa nešifrovali) a kódové označenia („čísla“ alebo jednoslabičné šifry), považované vtedy za spoľahlivú formu krytia správ. Kľúče k týmto šifram, často rozsiahle, starostlivo upratané tabuľky podľa abecedy nemožno rozoznávať rýchlo, sú to manuálne záznamy na papieri.

Manuálna analýza a rozdelenie odpovedajúcich tabuliek sa stáva tým náročná, chybováa vyžaduje veľa času. Každý kľúč sa líši svojim prevedením a kompozíciou (napr. tabuľka pre homofóne šifry, zoznam písmen pre slová a alebo n-grammi. Časti sú síce vizuálne oddelené, no ich presná identifikácia a miesto v rukou písaných kľúčoch predstavuje svoj druh výzvy. Táto úloha je sťažená rukopisom. Každého z rukopisov je iný a ľudí ktorí písali tieto kľúče pred stovkami rokov je úplne iný. 

Práve tu prichádza k slovu umelá inteligencia (AI), konkrétne metódy hlbokého učenia. Automatizované rozpoznávanie „nomenklátorových kľúčov“ z rukopisu by mohlo ušetriť veľa času a starosti histórom alebo kryptográfom. Identifikáciou a izoláciou jednotlivých funkčných blokov na stránke kľúča sa vytvára predpoklad pre ďalšie kroky analýzy a spracovanie. Tou je napríklad rozpoznávanie znakov (OCR) špecificky prispôsobené pre daný segmentovaný blok, alebo priama podpora pri dešifrovaní.

Cieľom práce je otestovanie a porovnani umelej inteligencie na automatickú identifikáciu kľúčových častí rukou písaných dokumentov. Tento cieľ sa dosiahne v rámci práce v týchto bodoch:   

1. Preštudovať a analyzovať problematiku historických substitučných šifier so zameraním na nomenklátory. Tento cieľ zahŕňa pochopenie ich štruktúry a vizuálneho usporiadania na rukou písaných dokumentoch.

2. Preskúmať súčasné metódy detekcie pomocou umelej inteligencie využiteľné pri rukou písaných dokumentov. Dôraz bude kladený na techniky objektovej detekcie, najmä na architektúry hlbokých neurónových sietí ako je YOLO. 

3. Analyzovať súčasný stav technológií a vybrať vhodný model umelej inteligencie pre úlohu segmentácie nomenklátorových kľúčov. Voľba bude zdôvodnená na základe relevantných kritérií. 

4. Pripraviť a anotovať dataset rukou písaných nomenklátorových šifrovacích kľúčov. Tento cieľ zahŕňa zber dát, definovanie anotácii a rozdelenie datasetu na tréningovú, validačnú a testovaciu množinu. Súčasťou bude aj príprava datasetu do troch foriem: pôvodnej, binarizovanej a augmentovanej. 

5. Natrénovať vybraný model umelej inteligencie na pripravených variantoch dátovej sady. Experimenty budú zahŕňať tréning na pôvodných, binarizovaných a augmentovaných dátach. 

6. Analyzovať a vyhodnotiť výsledky trénovania a testovania modelu. Porovnanie efektívnosti modelu na rôznych variantoch datasetu. Súčasťou bude aj testovanie na dátach, ktoré neboli použité pri tréningu, a návrh možných vylepšení implementovaného riešenia. 

7. Zdokumentovať celý proces, dosiahnuté výsledky a závery vo forme bakalárskej práce.

Naplnením týchto cieľov práca prispeje k automatizácii analýzy špecifického typu historických dokumentov pomocou moderných nástrojov umelej inteligencie.

Táto bakalárska práca je rozdelená do piatich kapitol, ktoré pokrývajú teóriu, analýzu, návrh riešenia, implementáciu a výsledky systému pre detekciu rukou písaných nomenklátorových šifrovacích kľúčov.

Kapitola Teoretické východiská poskytuje teoretický základ pre pochopenie problematiky. Najprv sa venuje historickým substitučným šifrám, s detailným zameraním na nomenklátory, ich typické časti a vizuálne charakteristiky na papieri. Následne sa kapitola zameriava na relevantné koncepty z oblasti umelej inteligencie pre porozumenie rukou písaných dokumentov, vrátane metód detekcie objektov s dôrazom na architektúru YOLO.

Kapitola Analýza súčasného stavu a návrh riešenia mapuje existujúce prístupy k problematike a prehľad využitia umelej inteligencie pri analýze historických dokumentov a šifier. Na základe tejto analýzy je zdôvodnený výber konkrétnej technológie pre riešenie problematiky a predstavený celkový návrh systému.

Kapitola Príprava dátovej sady (Datasetu) detailne popisuje proces tvorby a prípravy datasetu, ktorý je kľúčový pre tréning modelov umelej inteligencie. Zahŕňa informácie o zdroji a charakteristike dát, proces anotácie jednotlivých segmentov nomenklátorových kľúčov, definovaní tried, rozdelenie dát na tréningovú, validačnú a testovaciu množinu. Podrobne sú opísané aj kroky predspracovania dát a to príprava pôvodnej, binarizovanej a augmentovanej verzie datasetu.

Kapitola Implementácia a experimenty sa zaoberá realizáciou navrhnutého riešenia. Popisuje použité softvérové nástroje a knižnice, konfiguráciu zvoleného modelu YOLO (špecifikovanej verzie), nastavenie tréningových parametrov a samotný priebeh tréningových experimentov na troch pripravených variantoch dátovej sady. Záver kapitoly zahŕňa metriky použité na hodnotenie výkonnosti modelu.

Kapitola Výsledky a diskusia prezentuje výsledky dosiahnuté počas experimentačnej fázy. Zahŕňa vyhodnotenia výkonnosti modelu na pôvodnom, binarizovanom a augmentovanom datasete pomocou zvolených metrík. Výsledky sú porovnané, diskutuje sa vplyv predspracovania a augmentácie na presnosť segmentácie a je vykonaná kvalitatívna analýza úspešných aj problematických prípadov segmentácie. Kapitola tiež identifikuje limity navrhnutého riešenia a navrhuje možné zmeny pre jeho ďalšie vylepšenie.

Kapitola Záver sumarizuje zistenia a dosiahnuté výsledky práce vo vzťahu k stanoveným cieľom. Hodnotí prínos práce k automatizácii práce s rukou písanými dokumentami a navrhuje možné zmeny pre následné pokroky výskumu v danej oblasti.